% !TEX root = ../../hw05.tex
\begin{proof}[Solution]
  Following the hint, we first determine $f$ on the integers, then on
  the rationals, and finally on all of $\R$. Set $a=f(1)$.

  Taking $s=t=0$ in the additivity condition gives
  \[
    f(0)=f(0)+f(0),\qquad\text{so }f(0)=0.
  \]
  For any $t\in\R$, we also have
  \[
    0=f(t+(-t))=f(t)+f(-t),
  \]
  and therefore $f(-t)=-f(t)$. For a positive integer $m$, repeated
  addition gives
  \[
    f(m)=f(\underbrace{1+\cdots+1}_{m\text{ terms}})
      =mf(1)=ma.
  \]
  Together with $f(0)=0$ and $f(-m)=-f(m)$, this proves $f(m)=am$
  for every integer $m$.

  Now let $r=m/n$ be rational, where $m\in\Z$ and $n\in\N$.
  Adding $m/n$ to itself $n$ times gives
  \[
    nf\left(\frac mn\right)
      =f\left(\underbrace{\frac mn+\cdots+\frac mn}_{n\text{ terms}}\right)
      =f(m)=am.
  \]
  Dividing by $n$ yields
  \[
    f\left(\frac mn\right)=a\frac mn.
  \]
  Thus $f(r)=ar$ for every $r\in\Q$.

  \begin{pendingreview}
  Let $x\in\R$ be arbitrary.
  By the density of $\Q$ in $\R$, there exists a sequence $(q_n)$
  of rational numbers such that
  \[
    q_n\longrightarrow x.
  \]
  Since $f$ is continuous at $x$, the sequential characterization of
  continuity gives
  \[
    \lim_{n\to\infty}f(q_n)=f(x).
  \]
  But we have already established that $f(q_n)=aq_n$ for every $n$.
  Therefore,
  \[
    \begin{aligned}
      \lim_{n\to\infty}f(q_n)
        &=\lim_{n\to\infty}aq_n\\
        &=a\lim_{n\to\infty}q_n\\
        &=ax.
    \end{aligned}
  \]
  By the uniqueness of limits,
  \[
    f(x)=ax.
  \]
  Since $x\in\R$ was arbitrary, we conclude
  \[
    f(x)=ax,\qquad\forall x\in\R,\quad a=f(1).\qedhere
  \]
  \end{pendingreview}
\end{proof}

\begin{pendingreview}
The central insight is that additivity determines $f$ on all rational
numbers. Because the rationals are dense in the reals, continuity forces
the function to preserve that same linear relationship at every
irrational number.
\end{pendingreview}
